\chapter{Échanges avec les ingénieurs}
\label{chap:ingenieurs}

Le cœur de mon apprentissage a reposé sur des échanges avec trois ingénieurs occupant des fonctions de direction adjointe au sein de Reminex. Chacun m'a présenté un aspect différent du métier. Les développements ci-dessous restituent ce qui m'a été exposé lors de ces discussions~; il s'agit de sources orales, appuyées sur les notes que j'ai prises.

% ---------------------------------------------------------------------
\section{M. Sami Sayeh --- vision d'ensemble du Groupe}

M.~Sami~\textsc{Sayeh} m'a offert la première mise en perspective de mon stage. C'est lui qui m'a présenté le Groupe Managem et la chaîne de valeur minière détaillée au chapitre~\ref{chap:chaine}.

Au-delà de ce panorama, il a surtout mis en évidence les \textbf{liens entre trois grands mondes} qui doivent constamment dialoguer~: l'\textbf{exploration} (qui découvre et caractérise le gisement), l'\textbf{ingénierie} (qui conçoit l'usine et les infrastructures) et les \textbf{opérations} (qui exploitent). J'ai retenu que ces trois fonctions ne se succèdent pas simplement dans le temps~: elles s'informent mutuellement. L'ingénierie a besoin des données de l'exploration~; les opérations remontent vers l'ingénierie des retours de terrain qui nourrissent l'amélioration continue. Cette circulation d'informations entre métiers est une idée qui reviendra dans tout le rapport.

% ---------------------------------------------------------------------
\section{M. Houssam Benchakroun --- l'organisation de Reminex}

M.~Houssam~\textsc{Benchakroun} m'a présenté l'\textbf{organisation interne de Reminex}. À partir de ses explications et de mes notes, j'ai reconstitué le schéma suivant~:

\begin{center}
\begin{tikzpicture}[
  box/.style={rectangle, rounded corners=2pt, draw=darkslate!60, fill=darkslate!5,
    text width=3.1cm, minimum height=9mm, align=center, font=\small},
  head/.style={rectangle, rounded corners=2pt, draw=reminexorange!80, fill=reminexorange!12,
    text width=3.4cm, minimum height=9mm, align=center, font=\small\bfseries},
  sub/.style={rectangle, rounded corners=2pt, draw=centralegreen!60, fill=centralegreen!7,
    text width=3.4cm, minimum height=8mm, align=center, font=\footnotesize},
  arr/.style={-{Stealth[length=2.5mm]}, darkslate!60, line width=0.8pt}]
\node[head] (rem) {Reminex\\\scriptsize recherche minière \& exploration};
\node[box, below left=12mm and 6mm of rem] (ing) {Ingénierie};
\node[box, below=12mm of rem] (rd) {Centre R\&D};
\node[box, below right=12mm and 6mm of rem] (exp) {Exploration minière};
\draw[arr] (rem) -- (ing);
\draw[arr] (rem) -- (rd);
\draw[arr] (rem) -- (exp);
\node[sub, below=8mm of rd] (car) {\textcircled{1}~Caractérisation};
\node[sub, below=6mm of car] (pro) {\textcircled{2}~Procédé / \textit{flowsheet}};
\draw[arr] (rd) -- (car);
\draw[arr] (car) -- (pro);
\draw[arr, dashed] (rd.west) to[out=180,in=90] (ing.north);
\end{tikzpicture}
\captionof{figure}{Organisation interne de Reminex reconstituée d'après la présentation de M.~Benchakroun.}
\label{fig:org_reminex}
\end{center}

Reminex s'articule ainsi autour de trois pôles principaux~: l'\textbf{Ingénierie}, un \textbf{centre de Recherche \& Développement} et l'\textbf{Exploration minière}. M.~Benchakroun a particulièrement détaillé le rôle du centre de R\&D, qui réalise notamment~:

\begin{itemize}
\item la \textbf{caractérisation} du minerai (comprendre sa composition et son comportement)~;
\item le \textbf{développement du procédé et du \textit{flowsheet}} (définir comment le traiter).
\end{itemize}

Ce qui m'a paru important, c'est la \emph{collaboration} entre ces pôles~: le centre de R\&D fournit à l'Ingénierie le \textit{flowsheet} et les données que celle-ci va « industrialiser », tandis que l'Exploration alimente l'ensemble en connaissance du gisement. On retrouve, à l'échelle de Reminex, la même logique d'interdépendance que celle décrite par M.~Sayeh.

% ---------------------------------------------------------------------
\section{M. Yassine Ouhajjou --- l'ingénierie de procédé}

M.~Yassine~\textsc{Ouhajjou}, qui travaille en ingénierie de procédé, m'a fait la présentation la plus technique de mon stage. Son fil conducteur tenait en une phrase que j'ai notée telle quelle~: \textit{« transformer le \textit{flowsheet} en usine »}.

\subsection{Faire dialoguer toutes les disciplines}

Il a d'abord montré que passer d'un schéma de procédé à une usine réelle nécessite de mobiliser, autour du procédé, l'ensemble des disciplines~: génie civil, structure, mécanique, tuyauterie (\textit{piping}), instrumentation, ainsi que la démarche de sécurité intégrée dès la conception (\sig{sid}, \textit{Safety In Design}) et la prise en compte de l'environnement. Le procédé est au centre, mais il ne « tient debout » que grâce à toutes ces spécialités.

\subsection{Les documents clés de l'ingénierie}

M.~Ouhajjou a ensuite introduit les grands \textbf{livrables} qui décrivent progressivement l'usine, du plus général au plus détaillé~:

\begin{center}
\begin{tikzpicture}[
  doc/.style={rectangle, rounded corners=2pt, draw=reminexorange!80, fill=reminexorange!10,
    text width=2.5cm, minimum height=1.4cm, align=center, font=\footnotesize},
  arr/.style={-{Stealth[length=2.5mm]}, reminexorange!80, line width=1pt}]
\node[doc] (bfd) {\textbf{BFD}\\\scriptsize schéma par blocs};
\node[doc, right=8mm of bfd] (pfd) {\textbf{PFD}\\\scriptsize schéma de procédé};
\node[doc, right=8mm of pfd] (pid) {\textbf{P\&ID}\\\scriptsize tuyauterie \& instrumentation};
\node[doc, right=8mm of pid] (ds) {\textbf{Datasheets}\\\scriptsize fiches techniques équipements};
\draw[arr] (bfd) -- (pfd);
\draw[arr] (pfd) -- (pid);
\draw[arr] (pid) -- (ds);
\end{tikzpicture}
\captionof{figure}{Du plus général au plus détaillé~: les principaux livrables présentés par M.~Ouhajjou.}
\label{fig:documents}
\end{center}

Le \sig{bfd} donne la vue d'ensemble par grands blocs~; le \sig{pfd} précise les flux du procédé~; le \sig{pid} descend au niveau des tuyauteries et de l'instrumentation~; enfin les \textit{datasheets} spécifient chaque équipement. J'ai trouvé cette progression très éclairante~: elle illustre concrètement comment une idée abstraite devient un objet constructible.

\subsection{Le coût et la décision d'investir}

Un projet ne se juge pas seulement techniquement mais aussi \textbf{économiquement}. M.~Ouhajjou m'a présenté deux notions centrales~:

\begin{itemize}
\item le \sig{capex} (\textit{Capital Expenditure}), c'est-à-dire le coût d'investissement pour construire l'usine~;
\item l'\sig{opex} (\textit{Operating Expenditure}), c'est-à-dire le coût d'exploitation, une fois l'usine en marche (réactifs, énergie, main-d'œuvre\ldots).
\end{itemize}

La rentabilité se mesure ensuite par le retour sur investissement (\sig{roi}). Je précise que je n'ai réalisé \emph{aucun calcul} de ce type~: ces notions m'ont été expliquées pour que je comprenne la logique de décision, et non pour être mises en pratique.

\subsection{La maturation d'un projet par étapes}

Enfin, M.~Ouhajjou a insisté sur le fait qu'un projet minier ne se conçoit pas d'un seul coup~: il « mûrit » par étapes, chacune de plus en plus précise et de plus en plus coûteuse à produire. J'ai retenu l'enchaînement suivant~:

\begin{center}
\begin{tikzpicture}[
  ph/.style={rectangle, rounded corners=2pt, draw=centralegreen!70, fill=centralegreen!8,
    text width=2.55cm, minimum height=1.5cm, align=center, font=\scriptsize},
  arr/.style={-{Stealth[length=2.5mm]}, centralegreen!70, line width=1pt}]
\node[ph] (o) {\textbf{Étude\\d'orientation}\\\textit{(scoping)}\\précision indicative};
\node[ph, right=6mm of o] (pf) {\textbf{Pré-\\faisabilité}\\plusieurs scénarios};
\node[ph, right=6mm of pf] (f) {\textbf{Faisabilité}\\solution retenue\\($\approx\pm$\,30\,\%)};
\node[ph, right=6mm of f] (b) {\textbf{Ingénierie\\de base}\\dimensionnement};
\node[ph, right=6mm of b] (d) {\textbf{Ingénierie\\de détail}\\plans de construction};
\draw[arr] (o) -- (pf);
\draw[arr] (pf) -- (f);
\draw[arr] (f) -- (b);
\draw[arr] (b) -- (d);
\end{tikzpicture}
\captionof{figure}{Maturation d'un projet, de l'étude d'orientation à l'ingénierie de détail.}
\label{fig:maturation}
\end{center}

Ce séquencement m'a permis de comprendre une idée essentielle~: on n'investit massivement qu'après avoir progressivement réduit l'incertitude. La précision des estimations augmente à mesure qu'on avance dans les études, et c'est seulement à la fin que l'on engage la construction. C'est, en quelque sorte, le passage « du laboratoire à la production industrielle » que M.~Sayeh avait évoqué au début du stage.
